% =====================================================================
%  FICHE 8 — THÈME 3 — NIVEAU FACILE — CORRIGÉ
% =====================================================================
\documentclass[11pt,a4paper]{article}
\usepackage[utf8]{inputenc}
\usepackage[T1]{fontenc}
\usepackage{lmodern}
\usepackage[french]{babel}
\usepackage{amsmath,amssymb,mathtools}
\usepackage{icomma}
\usepackage{siunitx}
\sisetup{output-decimal-marker={,},per-mode=power,group-separator={\,},group-minimum-digits=5}
\usepackage[version=4]{mhchem}
\usepackage{xcolor}
\usepackage{tikz}
\usepackage[most]{tcolorbox}
\usepackage{enumitem}
\usepackage{array}
\usepackage{fancyhdr}
\usepackage[hidelinks]{hyperref}
\usetikzlibrary{arrows.meta,positioning,calc}
\usepackage[a4paper,margin=2.1cm,headheight=26pt,headsep=10pt]{geometry}
\usepackage{titlesec}

\definecolor{primary}{RGB}{150,98,162}
\definecolor{primarydark}{RGB}{92,52,110}
\definecolor{excol}{RGB}{0,135,90}

\tcbset{boxbase/.style={enhanced,breakable,boxrule=0.9pt,arc=2.5pt,left=3mm,right=3mm,top=2.3mm,bottom=2.3mm,fonttitle=\bfseries,coltitle=white,toptitle=0.6mm,bottomtitle=0.6mm}}
\newtcolorbox[auto counter]{correction}[1][]{boxbase,colback=excol!4,colframe=excol,title={\textsc{Corrigé}~\thetcbcounter\ \ifx\relax#1\relax\relax\else\textmd{--- \itshape #1}\fi}}

\newcommand{\rep}[1]{\textcolor{primarydark}{\textbf{#1}}}

\pagestyle{fancy}
\fancyhf{}
\fancyhead[L]{\small\textcolor{primarydark}{\textbf{Fiche 8} — Thème 3 : Calorimétrie}}
\fancyhead[R]{\small\textcolor{primarydark}{\textsc{Facile — corrigé}}}
\fancyfoot[C]{\small\thepage}
\renewcommand{\headrulewidth}{0.8pt}
\renewcommand{\headrule}{\hbox to\headwidth{\color{primary}\leaders\hrule height \headrulewidth\hfill}}
\setlength{\parindent}{0pt}\setlength{\parskip}{3pt}

\begin{document}

\begin{center}
\begin{tikzpicture}
\node[rectangle,rounded corners=7pt,fill=excol,text=white,minimum width=16cm,
      minimum height=1.1cm,align=center,inner sep=6pt]
  {{\large\bfseries Thème 3 — Calorimétrie}\\[1pt]{\small Niveau \textsc{Facile} — corrigé détaillé}};
\end{tikzpicture}
\end{center}

\begin{correction}[chaleur sensible]
\begin{enumerate}[label=\alph*),leftmargin=2em,topsep=2pt,itemsep=2pt]
  \item $\Delta T = 80 - 20 = \rep{\SI{60}{\kelvin}}$ (ou \SI{60}{\celsius}).
  \item $Q = m\,c\,\Delta T = 0{,}50 \times 4{,}185 \times 60 = \rep{\SI{125.55}{\kilo\joule}}$.
  \item L'eau \emph{se réchauffe} : elle \rep{reçoit} de la chaleur, donc $Q>0$.
\end{enumerate}
\end{correction}

\begin{correction}[chaleur latente]
\begin{enumerate}[label=\alph*),leftmargin=2em,topsep=2pt,itemsep=2pt]
  \item Fusion : $Q = m\,L_{\text{fus}} = 0{,}20 \times 334 = \rep{\SI{66.8}{\kilo\joule}}$.
  \item Vaporisation : $Q = m\,L_{\text{vap}} = 0{,}20 \times 2257 = \rep{\SI{451.4}{\kilo\joule}}$.
  \item La \rep{vaporisation} demande bien plus d'énergie : $\dfrac{2257}{334}\approx \rep{6{,}8}$ fois plus.
        (Vaporiser exige de séparer complètement les molécules, pas seulement de « décrocher » le réseau cristallin.)
\end{enumerate}
\end{correction}

\begin{correction}[sensible ou latente ?]
\begin{enumerate}[label=\alph*),leftmargin=2em,topsep=2pt,itemsep=2pt]
  \item Échauffement du liquide : \rep{$Q=mc\Delta T$}, $T$ \emph{varie}.
  \item Fusion : \rep{$Q=mL$}, $T$ \emph{constante} (\SI{0}{\celsius}).
  \item Échauffement de la vapeur : \rep{$Q=mc\Delta T$}, $T$ \emph{varie}.
  \item Vaporisation : \rep{$Q=mL$}, $T$ \emph{constante} (\SI{100}{\celsius}).
\end{enumerate}
\end{correction}

\begin{correction}[signe de la chaleur]
\begin{enumerate}[label=\alph*),leftmargin=2em,topsep=2pt,itemsep=2pt]
  \item Se réchauffe : $\Delta T>0$, donc \rep{$Q>0$} (chaleur reçue).
  \item Se refroidit : $\Delta T<0$, donc \rep{$Q<0$} (chaleur cédée).
  \item $Q_{\text{fer}} = 0{,}10 \times 0{,}45 \times (30-90) = 0{,}10 \times 0{,}45 \times (-60) =
        \rep{\SI{-2.7}{\kilo\joule}}$. $Q<0$ : le fer \rep{cède} de la chaleur (il refroidit).
\end{enumerate}
\end{correction}

\begin{correction}[la vigilance des unités]
\begin{enumerate}[label=\alph*),leftmargin=2em,topsep=2pt,itemsep=2pt]
  \item $\SI{3.5}{\kilo\joule} = \rep{\SI{3500}{\joule}}$.
  \item $Q = 0{,}10 \times 4{,}185 \times 10 = \rep{\SI{4.185}{\kilo\joule}} = \rep{\SI{4185}{\joule}}$.
  \item Il a oublié le facteur \rep{$1000$} entre \si{\kilo\joule} et \si{\joule} : $c$ étant en
        \si{\kilo\joule\per\kelvin\per\kilogram}, le résultat sort en \si{\kilo\joule}, soit
        $\SI{4.185}{\kilo\joule}=\SI{4185}{\joule}$ (et non \SI{4.185}{\joule}).
\end{enumerate}
\end{correction}

\end{document}
